\begin{enumerate}
\item L'opération qui permet de calculer 25\,\% de 480 est:

\begin{tabularx}{\linewidth}{*{4}{X}}
\textbf{a.~} $\dfrac{480}{25 \times 100}$&\textbf{b.~} $25 \times 480 \times 0,1$&
\textbf{c.~} $\dfrac{480 \times 100}{25}$&\textbf{d.~} $\dfrac 14 \times 480$
\end{tabularx}

\medskip

\begin{tabular}{@{\hspace{0.03\linewidth}} | p{0.95\linewidth}}
L'opération qui permet de calculer 25\,\% de 480 est
$480\times \dfrac{25}{100}$ soit $480\times \dfrac{1}{4}$.

\hfill\textbf{Réponse d.}
\end{tabular}

\bigskip

\item Voici trois nombres :

\begin{tabularx}{\linewidth}{*{3}{X}}
$A =\dfrac 15$& $B = \dfrac{19}{100}$& $C = 0,21$.
\end{tabularx}

Le classement par ordre croissant de ces trois nombres est :

\medskip

\begin{tabularx}{\linewidth}{*{4}{X}}
\textbf{a.~}$A < B < C$ &\textbf{b.~} A < C < B&\textbf{c.~} $B < A< C$&\textbf{d.~} $C < B < A$\\
\end{tabularx}

\medskip

\begin{tabular}{@{\hspace{0.03\linewidth}} | p{0.95\linewidth}}
$A=\dfrac{1}{5}=0,20$; $B=\dfrac{19}{100}=0,19$ et $C=0,21$

$0,19<0,20<0,21$ donc $B<A<C$
\hfill\textbf{Réponse c.}
\end{tabular}

\bigskip

\item Voici quatre nombres :

\begin{tabularx}{\linewidth}{*{4}{X}}
$A = \left(\dfrac 15\right)^2$ & $B = \left (\dfrac{1}{2}\right )^5$ & $C = 0,05$& $D = \left(\dfrac 13\right)^3$
\end{tabularx}

Le plus grand de ces quatre nombres est:

\medskip

\begin{tabularx}{\linewidth}{*{4}{X}}
\textbf{a.~} $A$&\textbf{b.~} $B$&\textbf{c.~} $C$&\textbf{d.~} $D$
\end{tabularx}

\medskip

\begin{tabular}{@{\hspace{0.03\linewidth}} | p{0.95\linewidth}}
$A=\left (\dfrac{1}{5}\right )^2= \dfrac{1}{5^2} =\dfrac{1}{25}$; $B=\left (\dfrac{1}{2}\right )^5=\dfrac{1}{2^5}=\dfrac{1}{32}$; $C=0,05 = \dfrac{5}{100}=\dfrac{1}{20}$ et

$D=\left (\dfrac{1}{3}\right )^3=\dfrac{1}{27}$
Toutes les fractions ont le même numérateur, donc la plus grande est celle qui a le plus petit dénominateur.
\hfill\textbf{Réponse c.}
\end{tabular}

\bigskip

\item Un article augmente de 10\,\% puis il augmente encore de 10\,\%.

Après ces deux augmentations il a augmenté de :

\begin{tabularx}{\linewidth}{*{4}{X}}
\textbf{a.~} $(10\,\%)^2$ &\textbf{b.~} 19\,\% &\textbf{c.~} 20\,\% &\textbf{d.~} 21\,\%
\end{tabularx}

\medskip

\begin{tabular}{@{\hspace{0.03\linewidth}} | p{0.95\linewidth}}
Augmenter de 10\,\%, c'est multiplier par $1+\dfrac{10}{100}=1,1$; augmenter deux fois de 10\,\%, c'est multiplier par $1,1\times 1,1=1,21$.

$1,21=1+\dfrac{21}{100}$ donc ça correspond à une augmentation de 21\,\%.

\hfill\textbf{Réponse d.}
\end{tabular}

\bigskip

\item Le tiers d'un quart correspond à la fraction :

\medskip

\begin{tabularx}{\linewidth}{*{4}{X}}
\textbf{a.~} $\dfrac 17$&\textbf{b.~} $\dfrac 34$& \textbf{c.~} $\dfrac 13 \times 4$& \textbf{d.~} $\dfrac{1}{12}$
\end{tabularx}

\medskip

\begin{tabular}{@{\hspace{0.03\linewidth}} | p{0.95\linewidth}}
Le tiers d'un quart correspond à la fraction $\dfrac{1}{3}\times \dfrac{1}{4}=\dfrac{1}{12}$.
\hfill\textbf{Réponse d.}
\end{tabular}

\bigskip

\item On considère $A = 10 + 0,1 + \dfrac{1}{1000}$. On a : 

\medskip

\begin{tabularx}{\linewidth}{*{4}{X}}
\textbf{a.~} $A= \dfrac{20^{-1}}{\np{1000}}$&\textbf{b.~} $A = \dfrac{1}{\np{1000}}$&\textbf{c.~} $A = 10,101$ &\textbf{d.~} $A = 10,110$ 
\end{tabularx}

\medskip

\begin{tabular}{@{\hspace{0.03\linewidth}} | p{0.95\linewidth}}
$A = 10 + 0,1 + \dfrac{1}{1000} = 10+0,1+0,001=10,101$
\hfill\textbf{Réponse c.}
\end{tabular}

\bigskip

\item On considère $A = 10^{10} + 10^{-10}$. $A$ est environ égal à : 

\begin{tabularx}{\linewidth}{*{4}{X}}
\textbf{a.~}$10^0$&\textbf{b.~} 0&\textbf{c.~} $10^{10}$&$100^0$
\end{tabularx}

\medskip

\begin{tabular}{@{\hspace{0.03\linewidth}} | p{0.95\linewidth}}
$10^{-10}$ es proche de 0 donc $A\approx 10^{10}$.
\hfill\textbf{Réponse c.}
\end{tabular}

\bigskip

\item Une durée de 100 minutes correspond à :

\begin{tabularx}{\linewidth}{*{4}{X}}
\textbf{a.~} 1 heure &\textbf{b.~} 1,40 heure &\textbf{c.~} $\dfrac 53$ heure&\textbf{d.~} 2 heures
\end{tabularx}

\medskip

\begin{tabular}{@{\hspace{0.03\linewidth}} | p{0.95\linewidth}}
$\dfrac{1}{3}$ heure correspond à 20 minutes, $\dfrac{2}{3}$ heure à 40 minutes.

$\dfrac{3}{3}$ heure correspondent à 60 minutes.
\hfill\textbf{Réponse c.}
\end{tabular}
\end{enumerate}

\medskip

\begin{minipage}{0.65\linewidth}
\begin{enumerate}[resume]
\item On considère une droite $D$ représentée ci-contre.

La seule équation pouvant correspondre à l'équation réduite
de la droite $D$ est :

\begin{tabularx}{\linewidth}{*{2}{X}}
\textbf{a.~} $y = x+3$ &\textbf{b.~} $y = x - 3$\\
\textbf{c.~} $y = -x + 3$ &\textbf{d.~} $y = - x - 3$
\end{tabularx}
\end{enumerate}
\end{minipage}\hfill
\begin{minipage}{0.33\linewidth}
\psset{xunit=0.5cm,yunit=0.3cm,arrowsize=2pt 3}
\begin{pspicture*}(-3,-4)(6,3)
\psaxes[linewidth=1.25pt,Dx=20,Dy=20]{->}(0,0)(-3,-4.5)(6,3)
\psplot[plotpoints=600,linewidth=1.25pt,linecolor=red]{-3}{6}{x 3 sub}
\uput[u](5.8,0){$x$} \uput[r](0,2.6){$y$}\uput[dr](5.8,2.8){\red $D$}
\end{pspicture*}
\end{minipage}

\medskip

\begin{enumerate}[label=]
\item \begin{tabular}{@{\hspace{0.03\linewidth}} | p{0.95\linewidth}}
La droite a un coefficient directeur positif et une ordonnée à l'origine négative.

\hfill\textbf{Réponse b.}
\end{tabular}
\end{enumerate}

\medskip

\begin{enumerate}[start=10]
\item On considère la fonction $f$ définie pour tout réel $x$ par $f(x) = 7 - \dfrac 12 (x - 3)^2$.

L'image de 3 par la fonction $f$ est égale à :

\begin{tabularx}{\linewidth}{*{4}{X}}
\textbf{a.~} $7 - \dfrac 12$ &\textbf{b.~} $7 - \dfrac 12(9 + 9)$ &\textbf{c.~} $7$&\textbf{d.~} 0
\end{tabularx}

\medskip

\begin{tabular}{@{\hspace{0.03\linewidth}} | p{0.95\linewidth}}
$f(3)=7-\dfrac{1}{2}\left (3-3\right )^2=7-0=7$
\hfill\textbf{Réponse c.}
\end{tabular}

\bigskip

\item Quand on développe $(x - 3)^2$ on obtient:

\begin{tabularx}{\linewidth}{*{4}{X}}
\textbf{a.~} $x^2 + 9$ &\textbf{b.~} $x^2 - 9$ &\textbf{c.~} $x^2 + 6x + 9$&\textbf{d.~} $x^2 - 6x + 9$
\end{tabularx}


\medskip

\begin{tabular}{@{\hspace{0.03\linewidth}} | p{0.95\linewidth}}
$(a-b)^2=a^2-2ab+b^2$
\hfill\textbf{Réponse d.}
\end{tabular}

\bigskip

\item Voici deux séries de valeurs :

Série A : 1 ; 2 ; 3 \qquad Série B : 0,5 ; 2 ; 100

Une seule de ces affirmations est exacte :

\begin{tabular}{l}
\textbf{a.~} Les deux séries ont la même moyenne et la même médiane.\\
\textbf{b.~} Les deux séries ont la même moyenne mais pas la même médiane.\\
\textbf{c.~} Les deux séries ont la même médiane mais pas la même moyenne.\\
\textbf{d.~} Les deux séries n'ont ni la même moyenne ni la même médiane.
\end{tabular}

\medskip

\begin{tabular}{@{\hspace{0.03\linewidth}} | p{0.95\linewidth}}
\begin{list}{\textbullet}{On regarde moyenne et médiane de chaque série.}
\item $1+2+3=6$ et $0,5+2+100=102,5$ donc les deux séries n'ont pas la même moyenne.
\item Les deux séries comportent 3 nombres donc leur médiane est le 2\ieme{} nombre, c'est-à-dire 2; les deux séries ont la même médiane.
\end{list}
\hfill\textbf{Réponse c.}
\end{tabular}


\end{enumerate}


\medskip
